% !PRINTED 61
\weberpage{61}
%<<CONT
\[
\begin{aligned}
A &\equiv -1 \pmod{2^{\lambda}} &&\equiv 1 \pmod{q_1^{\varkappa_1}} &&\equiv 1 \pmod{q_2^{\varkappa_2}} \ldots\\
B &\equiv \phantom{-}5\;'' &&\equiv 1\;'' &&\equiv 1\;'' \ldots\\
C_1 &\equiv \phantom{-}1\;'' &&\equiv g_1\;'' &&\equiv 1\;'' \ldots\\
C_2 &\equiv \phantom{-}1\;'' &&\equiv 1\;'' &&\equiv g_2\;'' \ldots
\end{aligned}
\]

\noindent
$\ldots\ldots\ldots\ldots\ldots\ldots\ldots\ldots\ldots\ldots\ldots\ldots\ldots\ldots\ldots\ldots\ldots\ldots\ldots\ldots\ldots\ldots$

and by 4.\ one then obtains, for every number $n$ of our group,

\weq{3}{n \equiv A^{\alpha} B^{\beta} C_1^{\gamma_1} C_2^{\gamma_2} \ldots \pmod{m}.}

The $A, B, C_1, C_2 \ldots$ are thus the elements of a basis of the
group $\mathfrak{N}$, of degrees $a, b, c_1, c_2 \ldots$

If $m$ is odd or divisible only by the first power of 2,
the two elements $A, B$ drop out of the basis. If $m$ is divisible
by 4, but not by 8, then $B$ drops out and the element $A$ of the
second degree remains.

\gesperrt{5. If one decomposes the numbers $a, b, c_1, c_2 \ldots$ into powers
of primes different from one another, one obtains the invariants of the group $\mathfrak{N}$.}

In order to comprehend the various cases together, we denote
the elements of the basis (2) by

\weq{4}{C_{-1}, C_0, C_1 \ldots C_\mu.}

Here $C_{-1}, C_0$ are to stand for $A$ and $B$, and are accordingly
to be set equal to 1 when $\lambda = 0$ or $\lambda = 1$. When $\lambda = 2$,
we are to set $B = C_{-1} = 1$, $A = C_0$, and when $\lambda > 1$, $A = C_{-1}$,
$B = C_0$.
% [sic: Weber prints \lambda > 1; the preceding cases require \lambda > 2 here.]

Let the degrees of these elements be denoted by

\weq{5}{c_{-1}, c_0, c_1 \ldots c_\mu}

and the indices of a number from $\mathfrak{N}$, which are to be taken with
the quantities (5) as moduli, by

\weq{6}{\nu_{-1}, \nu_0, \nu_1 \ldots \nu_\mu.}

If $\lambda = 0$ or 1, then $\nu_{-1}$ and $\nu_0$ have only the value 0;
if $\lambda = 2$, then $\nu_{-1}$ has the value 0 and $\nu_0$ the value 0 or 1.
If $\lambda \geqq 3$, then $\nu_{-1}$ has one of the two values 0, 1, and $\nu_0$
one of the values $0, 1, 2, \ldots c_0 - 1$; $\mu$ is always equal to the
number of odd primes, different from one another,
that divide $m$. Leaving $\nu_{-1}, c_{-1}$ aside, we shall
call $\nu_0, \nu_1 \ldots \nu_\mu$ the \gesperrt{indices} of $n$ corresponding to the primes $2, q_1 \ldots q_\mu$,
and $c_0, c_1 \ldots c_\mu$ the \gesperrt{index moduli} corresponding
to the same primes. These index moduli are
\[
c_0 = \varphi(2^{\lambda-1}), \quad c_1 = \varphi(q_1^{\varkappa_1}), \ldots c_\mu = \varphi(q_\mu^{\varkappa_\mu}),
\]
%CONT>>
